\section{Open problems}\label{sec:open}

\begin{enumerate}
\item \emph{Similar ordering through continued fractions.} Theorem~\ref{thm:similar} characterizes the similarly ordered Cohen--Macaulay four-element sets by the two Ap\'ery sets. Give a characterization by a short list of congruence or continued-fraction conditions on $(a_1,a_2,d)$, for instance through the Euclidean-type algorithm that computes the Ap\'ery sets of the three-generated semigroups $S_1$ and $S_2$ \cite{Rodseth1978}, and extend the diagram description and the test to five or more elements.
\item \emph{Other families.} The spectrum of $\{0,a,b,2b\}$ is governed by the three-distance theorem (Section~\ref{sec:family}). Compute the spectrum at the modulus $d$ of other complete-intersection families, such as $\{0,a,b,kb\}$ with $0<a<b$ coprime and $k\ge3$; for $k\ge3$, $b$ is not a modulus, since $b\notin S_2$.
\item \emph{Overshoot for more elements.} Every non-Cohen--Macaulay four-element set overshoots its mass at some level, but $\{0,5,9,11,20\}$ does not (Section~\ref{sec:noncm}). Which non-Cohen--Macaulay sets with five or more elements overshoot their mass? Proposition~\ref{prop:overshoot} and Theorem~\ref{thm:hinge} decide it for any given set by a finite computation; the question asks for a structural criterion. Hoa and Tien \cite[Thm.~4.4]{HoaTien2026} describe the points of $I_h$ outside $hA_0$, as the first local cohomology of $\Bbbk[\mathcal S]$, through Ap\'ery sets, which may be a starting point.
\end{enumerate}
